% Taksonomi CBDC berdasarkan segmen pengguna dan model distribusi (TikZ).
% Di-input dari BAB II §2.2.1.

\begin{figure}[H]
\centering
\resizebox{0.88\textwidth}{!}{%
\begin{tikzpicture}[
  root/.style={draw=black, fill=gray!30, rectangle, rounded corners=3pt,
    minimum width=2.8cm, minimum height=0.9cm, align=center, font=\small\bfseries},
  node1/.style={draw=black, fill=gray!12, rectangle, rounded corners=3pt,
    minimum width=3.0cm, minimum height=0.9cm, align=center, font=\footnotesize\bfseries},
  node2/.style={draw=black, fill=white, rectangle, rounded corners=3pt,
    minimum width=3.8cm, minimum height=1.6cm, align=center, font=\footnotesize},
  selected/.style={draw=black, fill=black!15, rectangle, rounded corners=3pt,
    minimum width=3.8cm, minimum height=1.6cm, align=center, font=\footnotesize\bfseries},
  ar/.style={-{Stealth}, thick},
]

%--- Akar ---
\node[root] (cbdc) at (5.5, 4.5) {CBDC};

%--- Level 1 ---
\node[node1] (wh) at (2.2, 3.0)  {Wholesale\\(Antarlembaga)};
\node[node1] (rt) at (8.8, 3.0)  {Retail\\(Masyarakat Umum)};

%--- Level 2 (Retail sub-types) ---
\node[node2] (one) at (6.4, 1.5)
  {Satu Jenjang\\(\textit{Direct})\\Bank Sentral\\$\to$ Publik};
\node[selected] (two) at (11.2, 1.5)
  {Dua Jenjang\\(\textit{Two-Tier})\\BI $\to$ Bank/PJP\\$\to$ Publik};

%--- Arrows ---
\draw[ar] (cbdc.south) to[out=-150,in=90] (wh.north);
\draw[ar] (cbdc.south) to[out=-30,in=90]  (rt.north);
\draw[ar] (rt.south)   to[out=-150,in=90] (one.north);
\draw[ar] (rt.south)   to[out=-30,in=90]  (two.north);

%--- Contoh ---
\node[font=\scriptsize, align=center] at (2.2, 1.8)
  {Contoh: Garuda PoC\\\cite{bi_garuda_poc}};
\node[font=\scriptsize, align=center, anchor=north] at (6.4, 0.45)
  {Contoh: DCash (Karibia)};
\node[font=\scriptsize, align=center, anchor=north] at (11.2, 0.45)
  {Contoh: eNaira, e-CNY,\\Rupiah Digital};

\end{tikzpicture}%
}
\caption{Taksonomi CBDC berdasarkan segmen pengguna dan model distribusi; model dua jenjang yang diadopsi pada penelitian ini ditampilkan dengan latar abu-abu sesuai rancangan Rupiah Digital Bank Indonesia}
\label{fig:cbdc-taxonomy}
\end{figure}
